\chapter{Literature Review}
\label{ch:literature}

\section{Overview}

This chapter traces the development of value-based reinforcement learning up to the two algorithms
compared here, and then examines what is and is not well established about the difference between
them. The organising question throughout is: \emph{whose} value is the agent learning?

\section{Markov decision processes and dynamic programming}

The formal setting is the Markov decision process, developed by Bellman \cite{bellman1957} and
given its modern treatment by Puterman \cite{puterman1994}. An MDP specifies states, actions,
transition probabilities and rewards, and the Markov property asserts that the next state depends
only on the current state and action.

Dynamic programming --- policy iteration and value iteration --- solves an MDP exactly, but
requires the transition and reward functions to be known in advance. A warehouse robot does not
have them. It has only the ability to act and observe what follows, which is the setting
reinforcement learning addresses.

\section{Monte Carlo and temporal-difference learning}

Two families of method learn from experience alone. Monte Carlo methods wait until an episode
terminates and update toward the observed return; they are unbiased but high-variance, and cannot
learn from incomplete episodes.

Temporal-difference learning, introduced by Sutton \cite{sutton1988}, instead updates toward a
\emph{bootstrapped} estimate --- the immediate reward plus the discounted value of the successor
state. TD methods learn online, from incomplete episodes, with far lower variance. The cost is
bias, since the target depends on the current estimate, which is itself wrong. Convergence results
for TD(0) under suitable step-size conditions are given by Dayan \cite{dayan1992} and
Jaakkola et al.\ \cite{jaakkola1994}.

The practical importance of this trade-off is that TD control learns \emph{while the robot is
still moving}, which is what makes the cost of exploration a measurable quantity in this study.

\section{Q-learning}

Q-learning, introduced by Watkins \cite{watkins1989} and proved convergent by Watkins and Dayan
\cite{watkins1992}, is the canonical \emph{off-policy} TD control algorithm. Its update
bootstraps from the maximum over next-state action values,
\[
Q(s,a) \leftarrow Q(s,a) + \alpha\Big[r + \gamma \max_{a'} Q(s',a') - Q(s,a)\Big],
\]
and this maximisation is what makes it off-policy: the target describes the greedy policy,
irrespective of the policy that generated the data. Given sufficient exploration and a suitable
step-size schedule, $Q$ converges to the optimal action-value function $Q^{*}$ regardless of how
the agent behaves while learning.

That independence is the algorithm's defining strength and, in the setting studied here, the source
of its characteristic weakness: because the target never refers to the behaviour policy, the cost
of the agent's own exploration never enters its estimates.

\section{SARSA}

SARSA was introduced by Rummery and Niranjan \cite{rummery1994} and named by Sutton
\cite{sutton1996} after the quintuple $(s,a,r,s',a')$ its update consumes:
\[
Q(s,a) \leftarrow Q(s,a) + \alpha\Big[r + \gamma Q(s',a') - Q(s,a)\Big],
\qquad a' \sim \pi(\cdot \mid s').
\]
Here $a'$ is the action the agent will actually take. The update therefore evaluates the behaviour
policy itself, making SARSA \emph{on-policy}. Under an $\varepsilon$-greedy policy, SARSA converges
to the optimal $\varepsilon$-greedy policy rather than to the optimal policy, and converges to
$Q^{*}$ only if $\varepsilon$ is decayed to zero appropriately --- the GLIE conditions
\cite{singh2000}.

This apparent weakness is the source of its practical strength. By evaluating the policy that
includes exploration, SARSA accounts for the possibility that the next action is random, and
therefore assigns lower value to states from which a random action would be catastrophic.

\section{The Cliff Walking result}

The clearest demonstration of the distinction is Example 6.6 of Sutton and Barto \cite{sutton2018},
\emph{Cliff Walking}. An agent crosses a grid from start to goal; a row of cliff cells lies between
them. Stepping into the cliff incurs a large penalty and returns the agent to the start.

The optimal path runs immediately alongside the cliff. Q-learning learns exactly that path, and
consequently falls off repeatedly while exploring, so its online performance during learning is
poor. SARSA learns a longer path further from the edge, because its estimates price in the
possibility of an $\varepsilon$-greedy step over it, and achieves better online return. Decaying
$\varepsilon$ makes both converge to the optimal policy.

Three features of this result are essential and are reproduced deliberately in the present study:

\begin{itemize}
\item The hazard is \textbf{localised and adjacent to the optimal path}. Without this there is no
      decision to make and the two algorithms behave identically.
\item The environment is \textbf{deterministic}. The effect is caused by $\varepsilon$-greedy
      exploration alone, not by stochastic dynamics.
\item The comparison measures \textbf{two distinct quantities} --- online return during learning,
      and the quality of the final greedy policy --- which have different answers.
\end{itemize}

The third point is the one most often lost in secondary accounts, which tend to report that
``SARSA is safer'' without noting that Q-learning's final policy is the better one.

\section{Related developments}

Expected SARSA \cite{vanseijen2009} replaces the sampled $Q(s',a')$ with its expectation under the
policy, reducing variance while remaining on-policy; it subsumes both algorithms as special cases.
Double Q-learning \cite{hasselt2010} addresses the maximisation bias that the $\max$ operator
introduces. Neither is evaluated here, since the object of study is the basic on-policy/off-policy
contrast, but both are natural extensions (Chapter~\ref{ch:conclusion}).

The distinction acquires additional importance under function approximation, where off-policy
bootstrapping combined with approximation and bootstrapping forms the ``deadly triad''
\cite{sutton2018} and can diverge outright. In the tabular setting studied here both algorithms
are convergent, so the difference observed is about the \emph{policy being evaluated} and not
about stability.

\section{Reinforcement learning for warehouse robots}

Automated guided vehicles and autonomous mobile robots in warehouses are a well-established
application area, and safety around human co-workers is a recognised constraint rather than an
afterthought. The field of safe reinforcement learning, surveyed by Garc\'ia and Fern\'andez
\cite{garcia2015}, is explicitly concerned with the problem this report measures: the cost incurred
\emph{during} learning, as distinct from the quality of the policy that results.

Within that framing, the on-policy/off-policy choice is one of the simplest available mechanisms
for trading final optimality against safety while learning, and it requires no additional
machinery --- only a different bootstrap target.

\section{Research gap}

The Cliff Walking result is textbook material, usually presented as a single illustrative run on a
$4 \times 12$ grid. What is less commonly available is that demonstration carried out with the
methodological care that would be expected of an empirical claim: multiple seeds, paired
statistical testing, an external optimality reference, a check that the effect generalises across
problem instances, and a sensitivity analysis over the hazard penalty that produces it.

This report supplies those, on a navigation problem whose obstacle layouts come from a published
dataset rather than from a hand-drawn figure, and it reports a negative control --- a floor without
a localised hazard, on which the celebrated distinction vanishes entirely.
